\documentclass[aspectratio=169]{beamer}
\usetheme{Madrid}
\usepackage{graphicx}
\usepackage{hyperref}
\usepackage{booktabs}

\title{Waymo: Autonomous Mobility \& Urban Impact}
\subtitle{Product Evaluation \& Policy Recommendations}
\author{Jiayi \and Daniel \and June}
\date{Urban Tech Seminar}

\begin{document}

\maketitle

\begin{frame}{Agenda}
\begin{itemize}
    \item \textbf{Product \& Context (Jiayi)}
    \begin{itemize}
        \item Technology \& Functionality
        \item Existing Deployments \& Market Trend
    \end{itemize}
    \vspace{0.2cm}
    \item \textbf{Impact Analysis (Daniel)}
    \begin{itemize}
        \item Physical Planning, Workforce Development, Data \& Privacy
    \end{itemize}
    \vspace{0.2cm}
    \item \textbf{Policy \& Design Recommendations (June)}
    \begin{itemize}
        \item Infrastructure Equity, AV-Transit Hubs, Grid Strategies
    \end{itemize}
\end{itemize}
\end{frame}

\section{Part 1: Product \& Context (Jiayi)}

\begin{frame}{What is Waymo?}
\textbf{Mission:} To make it safe and easy for people and things to get where they're going.
\vspace{0.5cm}

\textbf{Technology Stack (Level 4 Autonomy):}
\begin{itemize}
    \item \textbf{LiDAR:} 3D spatial mapping.
    \item \textbf{Radar:} Distance and speed tracking in all weather conditions.
    \item \textbf{Cameras:} Visual context (e.g., stoplights, signs).
    \item Relies heavily on high-fidelity, pre-mapped environments rather than vision alone.
\end{itemize}
\end{frame}

\begin{frame}{Current Deployments \& Market}
\begin{columns}[T]
\begin{column}{0.5\textwidth}
\textbf{Commercial Zones}
\begin{itemize}
    \item Phoenix (Waymo One)
    \item San Francisco
    \item Los Angeles
    \item Austin
\end{itemize}
\end{column}
\begin{column}{0.5\textwidth}
\textbf{Market Position}
\begin{itemize}
    \item Clear market leader in the US robotaxi sector.
    \item Transitioning from R\&D to scaled public utility.
\end{itemize}
\end{column}
\end{columns}
\end{frame}

\begin{frame}{User Experience \& Critiques}
The human-machine interaction introduces new friction points.

\vspace{0.5cm}
\textbf{The ``Empty Driver's Seat'' Dilemma}
\begin{itemize}
    \item Users often report feeling awkward or anxious without a human presence.
    \item \textbf{Proposal:} Implement a chatbot agent/avatar in the driver's seat screen to humanize the experience.
    \item \textbf{Edge Cases:} Managing passenger capacity limits (e.g., 5 people in a 4-person ride using the driver's seat) requires better warning systems and UI guidance during emergencies.
\end{itemize}
\end{frame}

\section{Part 2: Impact Analysis (Daniel)}

\begin{frame}{Physical Planning \& Infrastructure}
\begin{itemize}
    \item \textbf{Equity:} Waymo service zones often mirror historical disparities; better roads get AVs first, risking a ``two-tiered'' mobility system (UC Berkeley, 2025).
    \item \textbf{Energy Paradox:} ``Data centers on wheels.'' While the fleet avoids tailpipe emissions, compute power for autonomy requires massive grid energy (MIT, 2025).
    \item \textbf{Safety:} 92\% fewer pedestrian injuries compared to human drivers (Waymo, 2026).
    \item \textbf{Transit:} Potential to fill ``transit mobility gaps'' for seniors and paratransit riders (Valley Metro pilot).
\end{itemize}
\end{frame}

\begin{frame}{Workforce Development}
\begin{columns}[T]
\begin{column}{0.5\textwidth}
\textbf{Job Creation}
\begin{itemize}
    \item \textbf{+114k:} Projected new jobs in AV production, maintenance, and remote operations over 15 years (Chamber of Progress, 2024).
\end{itemize}
\end{column}
\begin{column}{0.5\textwidth}
\textbf{Labor Displacement}
\begin{itemize}
    \item 85\% of surveyed adults fear significant job losses for ride-hailing/delivery drivers (UCSD, 2026).
    \item Direct exposure to AV tech correlates with declines in commercial driver licensing.
\end{itemize}
\end{column}
\end{columns}
\end{frame}

\begin{frame}{Data \& Privacy}
\begin{itemize}
    \item \textbf{Surveillance Concerns:} Vehicles capture vast amounts of ``urban spatial data,'' raising concerns about increased street-level surveillance of neighborhoods.
    \item \textbf{Regulatory Confidentiality:} Waymo frequently requests keeping operational data confidential from regulators, citing trade secrets (SF Examiner).
    \item \textbf{Data Transparency:} Maintains an Open Dataset to share sensor data with researchers to improve safety perception models.
\end{itemize}
\end{frame}

\section{Part 3: Policy \& Design Recommendations (June)}

\begin{frame}{The Core Challenge}
\textbf{Domain Focus: Physical Planning}
\vspace{0.5cm}

How do we integrate AVs into the urban fabric to maximize safety and accessibility, without:
\begin{itemize}
    \item Exacerbating spatial inequities?
    \item Stressing the local power grid?
    \item Cannibalizing existing public transit networks?
\end{itemize}
\end{frame}

\begin{frame}{1. Equitable Service Zones}
\textbf{Mandate a ``Mobility Equity Framework'' for operators:}
\begin{itemize}
    \item Link expansion in high-income/high-demand areas with proportional service deployment in underserved neighborhoods.
    \item Condition operating permits on serving transit-challenged areas.
\end{itemize}
\vspace{0.5cm}
\textbf{Infrastructure Support:}
\begin{itemize}
    \item Municipalities should provide targeted infrastructure upgrades (e.g., road markings, signage) in lower-income areas to facilitate safe AV operation.
\end{itemize}
\end{frame}

\begin{frame}{2. AV-Transit Hubs}
\textbf{Design for First/Last Mile connectivity, not direct transit competition.}
\begin{itemize}
    \item \textbf{Physical Design:} Redesign infrastructure around major transit stations to include dedicated, geofenced AV pick-up/drop-off zones, repurposing curbside parking.
    \item \textbf{Pricing Policy:} Implement geofencing regulations that mandate lower pricing for AV trips that originate or terminate at high-capacity transit stations.
\end{itemize}
\end{frame}

\begin{frame}{3. Grid Integration \& V2I}
\textbf{Energy Management:}
\begin{itemize}
    \item Address the ``data center on wheels'' demand by zoning AV depots in industrial areas with robust electrical capacity.
    \item Mandate off-peak charging schedules or local renewable energy generation at fleet hubs.
\end{itemize}
\vspace{0.5cm}
\textbf{Vehicle-to-Infrastructure (V2I):}
\begin{itemize}
    \item Allocate infrastructure funds to upgrade traffic signals to communicate directly with the AV fleet.
    \item Standardize data sharing to continuously monitor intersection safety.
\end{itemize}
\end{frame}

\begin{frame}{Conclusion}
\begin{center}
\Large{\textbf{The Future of Urban Mobility}}
\vspace{0.5cm}

If regulated correctly, Waymo isn't just a technology product; it becomes a piece of public-private transit infrastructure that makes cities safer, more connected, and more accessible.
\vspace{1cm}

\textbf{Questions?}
\end{center}
\end{frame}

\end{document}
